Este reporte estudia el problema de hallar un camino 
de máximo puntaje en una grilla de $N \times M$ celdas, 
desde $(0,0)$ hasta $(N-1,M-1)$, moviéndose solo hacia 
abajo o a la derecha. Las celdas pueden ser libres 
(suman un valor), muros o minas 
(restan una penalización p y una vida); el camino es 
válido si llega a la meta con al menos una vida. 
Se implementaron tres algoritmos: fuerza bruta con 
backtracking, programación dinámica con estado 
(celda, minas pisadas) y un greedy que avanza hacia 
el vecino de mayor ganancia inmediata. El objetivo es 
contrastar su tiempo, memoria y calidad de solución con 
su complejidad teórica.